\section{A localized virial identity}
\label{sec:localization}

System Rellich--Pohozaev identities have a long history
\cite{Mitidieri1993,Souplet2009}; Phan gave a weighted version
\cite{Phan}. We use a potential-kernel form that retains the spatial
weights in a compactly supported interaction identity.

We retain the parameters and exact potentials of the preceding section.
Define the nonnegative interaction measure by
\begin{equation}\label{eq:interaction}
\begin{gathered}
  \dd\pi(x,y)=f(x)g(y)K(x-y)\dd x\dd y,\\
  \pi(\dd x,\R^n)=f(x)u(x)\dd x,
  \quad
  \pi(\R^n,\dd y)=g(y)v(y)\dd y.
\end{gathered}
\end{equation}
The marginal identities follow from Tonelli's Theorem and
\eqref{eq:potentials}. This measure is locally finite: on each compact
set the sources are integrable and the potentials are bounded.
Its total mass is not assumed finite. Values of the density or of the
quotients below on the diagonal may be set to zero.
For a unit vector \(e\in\Sph\), put
\[
  D_e(x,y)=\frac{((x-y)\cdot e)^2}{|x-y|^2},
  \qquad \omega_x=\frac{x}{|x|},
  \qquad D(x,y)=D_{\omega_x}(x,y).
\]
Choose any unit vector as \(\omega_0\); this convention does not affect
any integral.

The cutoff will vanish to a sufficiently high order at its outer boundary.
Writing \(t_+=\max\{t,0\}\), fix
\begin{equation}\label{eq:cutoff}
\begin{gathered}
  d(x)=\min\{1,(2-|x|)_+\},\qquad
  m=n+3,\qquad k=n(m+1),\qquad N=k+n+2,\\
  \eta(x)=\bigl[\bigl(1-(|x|-1)_+^3\bigr)_+\bigr]^N,
  \qquad w(x)=-x\cdot\nabla\eta(x),\qquad X(x)=x\eta(x).
\end{gathered}
\end{equation}
Thus \(\eta\in C_c^2(\R^n)\) equals one on \(B_1\), and \(w\ge0\)
vanishes unless \(1<|x|<2\). The cubic term makes the first two
derivatives match at radius one, and the power \(N\) makes them vanish
at radius two. On the intervening annulus, with \(s=|x|-1\),
\(1-s^3=d(x)(1+s+s^2)\). Differentiation therefore gives
\begin{equation}\label{eq:cutoff-bounds}
\begin{gathered}
  d^N\le\eta\le C d^N,\qquad
  |\nabla\eta|+w\le C d^{N-1},\\
  |\nabla^2\eta|+|\nabla w|+|\nabla^2X|\le C d^{N-2},
  \qquad dw\le C\eta.
\end{gathered}
\end{equation}
Set
\[
\begin{gathered}
  E_f=\int_{\R^n}\eta fu\dd x,\qquad
  E_g=\int_{\R^n}\eta gv\dd x,\qquad E=E_f+E_g,
  \\
  W_f=\int_{\R^n}wfu\dd x,\qquad
  W_g=\int_{\R^n}wgv\dd x.
\end{gathered}
\]
All these quantities are finite for each solution: \(u,v\) are bounded
on \(\overline B_2\), and \(|x|^A,|x|^B\) are integrable there because
\(A,B>-2\) and \(n\ge3\). This asserts finiteness only; a bound
independent of the solution remains to be proved.

\begin{lemma}[Localized virial identity]
The following identity holds, with an absolutely convergent integral:
\begin{equation}\label{eq:virial}
\begin{split}
  a(n+B)E_f+b(n+A)E_g-aW_f-bW_g
    &=(n-2)\int_{\R^{2n}}\Phi(x,y)\dd\pi(x,y),\\
  \Phi(x,y)&=
    \frac{(X(x)-X(y))\cdot(x-y)}{|x-y|^2}.
\end{split}
\end{equation}
\end{lemma}

\begin{proof}
We first justify absolute convergence of the separate kernel-gradient
integrals, before combining them. For \(r=|x|\) sufficiently small,
local boundedness of \(v\) gives
\(g(y)\le C_{u,v}r^A\) when \(|x-y|\le r/2\). Hence
\[
  |X(x)|\int_{|x-y|\le r/2}|x-y|^{1-n}g(y)\dd y
       \le C_{u,v}r\,r^A\int_0^{r/2}\dd s
       \le C_{u,v}r^{2+A}.
\]
On the complementary region, \(|X(x)|\le r\le2|x-y|\), so the same
expression with that region in place of the ball is at most
\(C(K*g)(x)=Cu(x)\). Both bounds remain bounded as \(x\to0\).
On compact subsets away from zero, the nearby source is bounded and
\(|\nabla K|\) is locally integrable; the remaining part is bounded
by a constant times the finite potential. Since \(X\) is supported
in \(\overline B_2\) and \(f\) is locally integrable, this proves
\[
  \int_{\R^n}f(x)|X(x)|
       \int_{\R^n}|\nabla K(x-y)|g(y)\dd y\dd x<\infty.
\]
Interchanging \((f,u,B)\) and \((g,v,A)\) proves the other assertion.
In particular, \(fX\cdot\nabla u\) and \(gX\cdot\nabla v\) are
absolutely integrable.

Since \(a=1/(q+1)\), integration by parts outside \(B_\rho\),
with \(0<\rho<1\), gives
\[
\begin{split}
  -\int_{|x|>\rho} fX\cdot\nabla u\dd x
   ={}&a\int_{|x|>\rho}|x|^B u^{q+1}
             \bigl((n+B)\eta-w\bigr)\dd x\\
     &+a\rho^{B+1}\int_{\partial B_\rho}u^{q+1}\dd S.
\end{split}
\]
Here
\(\operatorname{div}(|x|^BX)=|x|^B((n+B)\eta-w)\), and the
boundary term has the displayed positive sign because the outward normal
of the punctured domain on \(\partial B_\rho\) is \(-\omega_x\).
Its magnitude is \(O_{u,v}(\rho^{n+B})\), which tends to zero.
The corresponding term for \(gX\cdot\nabla v\) is positive and
\(O_{u,v}(\rho^{n+A})\), and also tends to zero. There is no outer
boundary term because \(X\) is compactly supported. Consequently,
\[
  -\int fX\cdot\nabla u\dd x-\int gX\cdot\nabla v\dd x
       =a(n+B)E_f+b(n+A)E_g-aW_f-bW_g.
\]
Finally substitute the differentiated potential identities and
\(\nabla K(z)=-(n-2)zK(z)/|z|^2\). Absolute convergence permits
Fubini's Theorem and interchanging the variables in the second integral;
their sum is exactly the right side of \eqref{eq:virial}.
\end{proof}

\begin{lemma}[Localization]\label{lem:localization}
For every \(\eps>0\), there is a constant \(C_\eps\), independent of
the solution, such that
\begin{equation}\label{eq:localization}
  |E_f-E_g|+
  \int_{\R^{2n}}|\Phi(x,y)-\eta(x)+w(x)D(x,y)|\dd\pi(x,y)
       \le\eps E+C_\eps.
\end{equation}
\end{lemma}

\begin{proof}
Fix \(0<\delta<1/4\), and write \(\ell(z)=\delta d(z)^m\) whenever
\(d(z)>0\). Call a pair \((x,y)\) close if
\(|x-y|\le\ell(z)\) for at least one endpoint \(z\in\{x,y\}\)
with \(d(z)>0\), and far otherwise. Since \(d\) is 1-Lipschitz,
for a close pair every point \(\xi\) of the segment \([x,y]\) satisfies
\[
  (1-\delta)d(z)\le d(\xi)\le(1+\delta)d(z).
\]
In particular, the entire segment is inside \(B_2\).
The derivative matrix of \(X\) is
\[
  \nabla X(x)=\eta(x)\operatorname{Id}
                   -w(x)\omega_x\otimes\omega_x,
\]
with its second term zero on \(B_1\). Taylor's formula and
\eqref{eq:cutoff-bounds} give, on close pairs with \(x\ne y\),
\begin{align*}
  |\Phi-\eta(x)+w(x)D|
       &\le C\delta d(z)^{N-2+m}
        =C\delta d(z)^{N+n+1},\\
  |\eta(x)-\eta(y)|
       &\le C\delta d(z)^{N-1+m}
        =C\delta d(z)^{N+n+2}.
\end{align*}
Both are at most \(C\delta d(z)^N\). Since
\(d(z)^N\le\eta(z)\le\eta(x)+\eta(y)\), the marginal identities
in \eqref{eq:interaction} bound the integral of their sum over close
pairs by \(C\delta E\).

For a far pair, every endpoint \(z\) in \(B_2\) satisfies
\(|x-y|>\ell(z)\). Using
\(|\Phi|\le(|X(x)|+|X(y)|)/|x-y|\) and \(0\le D\le1\), we obtain
\[
  |\Phi-\eta(x)+w(x)D|+|\eta(x)-\eta(y)|
   \le C\sum_{\substack{z\in\{x,y\}\\d(z)>0}}
          \left(d(z)^{N-1}+\frac{d(z)^N}{\ell(z)}\right).
\]
Pairs with both endpoints outside \(B_2\) contribute zero.
For the terms associated with \(x\), the tail estimate in
Proposition~\ref{prop:uniform}, \eqref{eq:uniform}, bounds the integral
by
\[
\begin{split}
  &C\int_{B_2}f(x)
       \left(d(x)^{N-1}+\frac{d(x)^N}{\ell(x)}\right)
       \ell(x)^{2-n}\dd x\\
  &\quad\le C\int_{B_2}f(x)
      \left(\delta^{2-n}d(x)^{4n+7}
                 +\delta^{1-n}d(x)^{3n+5}\right)\dd x
       \le C_\delta.
\end{split}
\]
Indeed the two exponents are
\[
  N-1-m(n-2)=4n+7>0,\qquad
  N-m(n-1)=3n+5>0,
\]
and \(\int_{B_2}f\le C\) by \eqref{eq:uniform}.
The terms associated with \(y\) satisfy the same bound with \(g\).
All these constants are independent of the solution.

Finally, both cutoff marginals are finite, so
\[
  |E_f-E_g|
    =\left|\int_{\R^{2n}}(\eta(x)-\eta(y))\dd\pi(x,y)\right|
    \le\int_{\R^{2n}}|\eta(x)-\eta(y)|\dd\pi(x,y).
\]
The preceding estimates bound the left side of
\eqref{eq:localization} by \(C\delta E+C_\delta\).
Choosing \(\delta\) sufficiently small proves the result.
\end{proof}

To obtain a uniform bound for \(E\), it remains to prove that, for every
\(\eps>0\),
\[
 aW_f+bW_g-(n-2)\int w(x)D(x,y)\dd\pi(x,y)
 \le\eps E+C_\eps.
\]
The next two sections establish this comparison from source-layer
intervals. In Section~\ref{sec:conclusion}, it will combine with the
marginal estimate above and the positive subcritical gap \(\gamma\)
to give the energy bound.
